\documentclass[12pt, a4paper]{report}

% ─── Encodage & langue ───────────────────────────────────────────────────────
\usepackage[utf8]{inputenc}
\usepackage[T1]{fontenc}
\usepackage[french]{babel}

% ─── Mise en page ────────────────────────────────────────────────────────────
\usepackage[top=2.5cm, bottom=2.5cm, left=2.5cm, right=2.5cm]{geometry}
\usepackage{setspace}
% Interligne : 1.15 est un bon compromis lisibilité / volume quand l'école impose
% une fourchette de pages. Si ton école impose un interligne, respecte-le.
% ⚠️ Piège : \onehalfspacing vaut déjà ~1,25 à 12 pt — passer de l'un à 1,25 ne gagne rien.
\setstretch{1.15}

% ─── Typographie ─────────────────────────────────────────────────────────────
\usepackage{lmodern}
\usepackage{microtype}
% Marge de manœuvre pour les paragraphes contenant de longs blocs \texttt{} insécables
% (noms de variables, routes d'API), qui déborderaient sinon dans la marge.
\setlength{\emergencystretch}{3em}

% ─── Graphiques & couleurs ───────────────────────────────────────────────────
\usepackage{graphicx}
\usepackage[dvipsnames, table]{xcolor}
\usepackage{tikz}
\usetikzlibrary{arrows.meta, positioning, backgrounds, fit, calc}
\usepackage{float}

% ─── Tableaux ────────────────────────────────────────────────────────────────
\usepackage{booktabs}
\usepackage{tabularx}
\usepackage{multirow}

% ─── Code source ─────────────────────────────────────────────────────────────
\usepackage{listings}
\usepackage{lstautogobble}

\definecolor{codebg}{RGB}{245,245,245}
\definecolor{codeframe}{RGB}{200,200,200}
\definecolor{codecomment}{RGB}{100,100,100}
\definecolor{codekeyword}{RGB}{0,0,180}
\definecolor{codestring}{RGB}{163,21,21}

\lstset{
  backgroundcolor=\color{codebg},
  basicstyle=\ttfamily\small,
  breakatwhitespace=false,
  breaklines=true,
  captionpos=b,
  commentstyle=\color{codecomment}\itshape,
  frame=single,
  rulecolor=\color{codeframe},
  keepspaces=true,
  keywordstyle=\color{codekeyword}\bfseries,
  numbers=left,
  numbersep=8pt,
  numberstyle=\tiny\color{gray},
  showstringspaces=false,
  stringstyle=\color{codestring},
  tabsize=2,
  autogobble=true
}

% ─── Liens & PDF ─────────────────────────────────────────────────────────────
\usepackage{amsmath}
\usepackage[colorlinks=true, linkcolor=NavyBlue, citecolor=NavyBlue, urlcolor=NavyBlue]{hyperref}
\usepackage{bookmark}
\usepackage[french]{cleveref}
% Fix INDISPENSABLE : babel français rend ":" actif, ce qui casse \cref{fig:xxx}
\AtBeginDocument{\shorthandoff{:}}

% ─── Titres de chapitre ──────────────────────────────────────────────────────
% La classe report réserve ~5 cm de blanc avant chaque titre de chapitre ; on ramène
% cet espace à une valeur raisonnable (utile quand le nombre de pages est contraint).
\usepackage{titlesec}
\titlespacing*{\chapter}{0pt}{0pt}{25pt}

% ─── En-têtes & pieds de page ────────────────────────────────────────────────
\usepackage{fancyhdr}
\setlength{\headheight}{14.5pt}
\addtolength{\topmargin}{-2.5pt}
\pagestyle{fancy}
\fancyhf{}
\fancyhead[L]{\small\leftmark}
\fancyhead[R]{\small <<PRÉNOM NOM>>}
\fancyfoot[C]{\thepage}
\renewcommand{\headrulewidth}{0.4pt}

% ─── Césure de noms propres récalcitrants ────────────────────────────────────
% Ajoute ici les mots que LaTeX coupe mal (noms de villes, d'organismes…).
\hyphenation{}

% ─── Numérotation ────────────────────────────────────────────────────────────
\setcounter{secnumdepth}{3}
\setcounter{tocdepth}{3}

% ─── Bibliographie ───────────────────────────────────────────────────────────
\usepackage{csquotes}
\usepackage[backend=biber, style=ieee, sorting=none]{biblatex}
\addbibresource{references.bib}
\setlength{\bibitemsep}{0.4\baselineskip}

% ─── Glossaire ───────────────────────────────────────────────────────────────
\usepackage[acronym, toc]{glossaries}
\makeglossaries
\newacronym{api}{API}{Application Programming Interface}
\newacronym{ci}{CI}{Continuous Integration}
\newacronym{hmoe}{HMoE}{Heterogeneous Mixture of Experts}
\newacronym{iot}{IoT}{Internet of Things}
\newacronym{json}{JSON}{JavaScript Object Notation}
\newacronym{llm}{LLM}{Large Language Model}
\newacronym{ot}{OT}{Operational Technology}
\newacronym{peft}{PEFT}{Parameter-Efficient Fine-Tuning}
\newacronym{qlora}{QLoRA}{Quantized Low-Rank Adaptation}
\newacronym{sse}{SSE}{Server-Sent Events}

\newglossaryentry{graphedattaque}{
  name={graphe d'attaque},
  description={Représentation des dépendances et chemins permettant de relier
  plusieurs étapes d'une compromission}
}

\newglossaryentry{veriteterrain}{
  name={vérité terrain},
  description={Ensemble de résultats attendus utilisé comme référence lors de
  l'évaluation d'un système}
}


% ─── Métadonnées PDF ─────────────────────────────────────────────────────────
\hypersetup{
  pdftitle={<<TITRE DU RAPPORT>>},
  pdfauthor={<<PRÉNOM NOM>>},
  pdfsubject={Rapport de stage — <<ÉCOLE>>},
}

% ─────────────────────────────────────────────────────────────────────────────
\begin{document}

% ─── Page de garde ───────────────────────────────────────────────────────────
\begin{titlepage}
  \centering

  % ─── Logos ───────────────────────────────────────────────────────────────
  % Dépose les fichiers dans figures/. La plupart des écoles exigent le logo
  % de l'école ET celui de l'organisme d'accueil.
  \begin{minipage}{0.45\textwidth}
    \centering
    % Remplace le cadre par : \includegraphics[height=2cm]{figures/logo-ecole.png}
    \framebox[4cm][c]{\rule{0pt}{2cm}\small logo école}\\
    {\Large\textbf{TELECOM Nancy}}\\
    {\small 2A}
  \end{minipage}
  \hfill
  \begin{minipage}{0.45\textwidth}
    \centering
    % Remplace le cadre par : \includegraphics[height=2cm]{figures/logo-organisme.png}
    \framebox[4cm][c]{\rule{0pt}{2cm}\small logo organisme}\\[0.3em]
    {\small <<SERVICE / DÉPARTEMENT — ORGANISME>>}
  \end{minipage}

  \vspace{1.5cm}
  \rule{\textwidth}{0.5pt}
  \vspace{0.5cm}

  {\small <<MENTION DU SEMESTRE / TYPE DE STAGE>>}

  \vspace{1cm}

  {\LARGE\bfseries <<TITRE DU RAPPORT\\[0.3em]
  sur deux ou trois lignes>>}

  \vspace{0.5cm}
  \rule{\textwidth}{0.5pt}

  \vspace{2cm}

  \begin{minipage}{0.48\textwidth}
    \textbf{Stagiaire}\\
    Léo Basin\\
    TELECOM Nancy, 2A
  \end{minipage}
  \hfill
  \begin{minipage}{0.48\textwidth}
    \textbf{Maître de stage}\\
    <<PRÉNOM NOM>>\\
    <<SERVICE>>\\
    <<ORGANISME>>
  \end{minipage}

  \vspace{1.5cm}

  \begin{minipage}{0.48\textwidth}
    \textbf{Tuteur école}\\
    <<PRÉNOM NOM>>\\
    TELECOM Nancy
  \end{minipage}

  \vspace{2cm}

  {\large Période du stage}\\[0.3em]
  {\large <<DATE DÉBUT>> — <<DATE FIN>>}

  \vfill

  {\small <<VILLE, PAYS>> \quad|\quad Année universitaire <<20XX--20XX>>}
\end{titlepage}

\clearpage

% ─── Remerciements ───────────────────────────────────────────────────────────
\chapter*{Remerciements}
\addcontentsline{toc}{chapter}{Remerciements}
\textbf Je tiens à remercier mes encadrants pour ce stage, M. Saïd Mahmoudi ainsi que M. Thibault Cholez. J’étends ces remerciements à M. Tanguy Vansnick et à M. Gaspard Menou qui m’ont aidé pour l’utilisation de l’outil LANCE et mon adaptation à celui-ci.


% ─── Résumés ─────────────────────────────────────────────────────────────────
\chapter*{Résumé}
\addcontentsline{toc}{chapter}{Résumé}
LANCE est un agent fondé sur un modèle de langage pour conduire et évaluer des
audits de réseaux IoT/OT. La mission principale du stage a consisté à concevoir
une architecture HMoE agentique~: un modèle de base quantifié est partagé entre
des adaptateurs QLoRA spécialisés pour l'organisation, la reconnaissance,
l'analyse de vulnérabilités et l'exploitation. Une chaîne de préparation
transforme les traces en exemples validés, tandis qu'un routeur déterministe
sélectionne l'expert associé à chaque phase derrière une API commune. Le
pipeline encadre les transitions, les outils, les sorties et les preuves. En
parallèle, le harness expérimental déploie des scénarios reproductibles et mesure
le comportement observé pendant les campagnes. L'évaluation associe des
métriques déterministes à un jugement sémantique encadré, sans confier le calcul
final au modèle. Des mécanismes de maîtrise du contexte complètent
l'architecture. Les premiers résultats
mettent en évidence un écart entre constats produits et chemins d'attaque
entièrement prouvés. Ils doivent encore être confirmés par des campagnes répétées
comparant les experts spécialisés à une référence généraliste dans des
conditions identiques.

\vspace{1em}
\noindent\textbf{Mots-clés~:} agent LLM, cybersécurité IoT/OT, HMoE.


% ─── Glossaire ───────────────────────────────────────────────────────────────
% Pas de \glsaddall : seules les entrées réellement employées dans le texte sont
% imprimées, ce qui évite un glossaire rempli d'entrées jamais citées.
% Appareil de référence composé en interligne simple (usage courant).
\begin{singlespace}
\printglossary[type=\acronymtype, title=Liste des acronymes]
\printglossary[title=Glossaire]
\end{singlespace}
\clearpage

% ─── Table des matières ──────────────────────────────────────────────────────
\begin{singlespace}
\tableofcontents
\end{singlespace}
\clearpage

% ─── Corps du rapport ────────────────────────────────────────────────────────
\chapter{Introduction}
\label{chap:introduction}

Les objets connectés et les systèmes industriels forment des réseaux
hétérogènes dans lesquels une passerelle, un courtier MQTT, un automate ou un
service web peut devenir une étape d'un chemin d'attaque. Évaluer leur sécurité
suppose de découvrir la topologie, de choisir des actions adaptées aux
protocoles, puis de vérifier les hypothèses par des observations et d'en
conserver les preuves dans des artefacts comparables entre plusieurs campagnes.
Les graphes d'attaque offrent un cadre pour représenter les dépendances entre
connectivité, configuration et compromission\cite{ou2005mulval}.

Le projet LANCE (\emph{LLM Agent for Network Compromise Evaluation}) étudie
l'emploi de \glspl{llm} pour automatiser cette démarche dans des réseaux
\gls{iot}/\gls{ot}. Son agent conduit une campagne structurée, depuis la
découverte du réseau jusqu'au rapport. Cette application pose une difficulté
centrale~: un modèle généraliste peut produire une réponse plausible sans
respecter le rôle attendu, le format d'un appel d'outil ou le niveau de preuve
nécessaire à une conclusion de sécurité.

La mission principale du stage a consisté à concevoir une architecture agentique
de type \gls{hmoe}, fondée sur un modèle de base partagé et sur des experts
spécialisés par \gls{qlora}. Chaque expert apprend un rôle opérationnel
organisation de la campagne, reconnaissance, analyse de vulnérabilités ou
exploitation et un routeur sélectionne l'adaptateur correspondant à la phase
exécutée. Le travail couvre ainsi la préparation de données issues de traces,
leur contrôle, l'entraînement, le service des adaptateurs et leur
intégration au cycle d'action de l'agent.

Le service HMoE est déployé localement derrière une API compatible OpenAI. Le
pipeline peut ainsi l'utiliser comme un fournisseur de modèle sans exposer la
sélection interne des experts.

Le pipeline d'exécution et le harness expérimental constituent le support de
cette spécialisation. Le premier impose les transitions, les outils autorisés
et les contrats de sortie ; le second déploie des scénarios reproductibles,
conserve la vérité terrain du côté de l'évaluateur et mesure le comportement
observé pendant une campagne. Leur consolidation a porté sur la robustesse de
l'orchestration, la collecte de preuves, la gestion des exécutions et
l'évaluation déterministe et sémantique.

La problématique peut dès lors être formulée ainsi~:
\emph{comment spécialiser un agent fondé sur un LLM local pour les différentes
étapes d'un audit IoT/OT, puis démontrer la qualité de son comportement au moyen
d'un protocole d'exécution et d'évaluation reproductible ?}

Le \cref{chap:contexte} présente le cadre et les objectifs de la mission. Le
\cref{chap:implementation} décrit l'intégration de la solution, notamment le
pipeline et le harness. Enfin, le \cref{chap:science} expose \gls{qlora}, le
routage \gls{hmoe} et leur évaluation.

\chapter{Contexte et présentation de la mission}
\chaptermark{Contexte et présentation}
\label{chap:contexte}

\section{Cadre du projet}

\subsection{Organisme d'accueil}

Le stage s'est déroulé à l'Université de Mons (UMONS), au sein du service
Informatique, Logiciel et Intelligence Artificielle (ILIA) de la Faculté
Polytechnique, implanté rue de Houdain à Mons. Ce service mène des activités
d'enseignement et de recherche en informatique. Ses travaux couvrent notamment
l'intelligence artificielle, la vision par ordinateur, la gestion de données
issues de l'Internet des objets ainsi que les environnements Cloud et
Edge~Computing. Ses enseignements incluent également le génie logiciel, la
programmation embarquée, l'IoT et la cybersécurité.

L'encadrement du stage a été assuré par M.~Saïd Mahmoudi, membre du service
ILIA. Ce cadre associait ainsi les compétences du service
en intelligence artificielle, systèmes distribués et IoT aux problématiques de
sécurité traitées par LANCE.

\subsection{Organisation technique}

Le développement s'appuie sur deux espaces aux responsabilités distinctes. Le
dépôt canonique rassemble le code versionné de l'agent, de l'\gls{api}, de
l'interface, du benchmark et de l'infrastructure. Il constitue la référence pour
les fonctionnalités, les configurations, les tests et la documentation associés
à une version donnée. L'espace d'entraînement accueille les traces préparées,
les jeux de données, les checkpoints et les adaptateurs produits pendant les
expériences.

Cette séparation évite de placer dans le dépôt des artefacts volumineux ou des
données sensibles, tout en maintenant la traçabilité expérimentale. Une
campagne doit pouvoir être reliée à une configuration du modèle, à une version
du pipeline, à un jeu de données identifié et à un protocole d'évaluation.
Les scripts et configurations restent donc versionnés, tandis que données et
poids sont conservés dans l'espace prévu à cet effet.

\section{Système existant}

LANCE associe trois ensembles fonctionnels~:

\begin{description}
  \item[Agent d'audit] L'agent conduit une campagne en six phases, de l'analyse
  initiale du réseau à la synthèse. Chaque phase reçoit un contexte défini,
  dispose d'outils autorisés et produit un artefact structuré pour la suivante.

  \item[Pipeline d'exécution] Le pipeline organise les transitions, parallélise
  certaines analyses, valide les sorties et applique les règles de sécurité.
  Il transforme les décisions probabilistes en actions observables et
  contrôlables.

  \item[Harness expérimental] Le benchmark décrit des réseaux
  \gls{iot}/\gls{ot} reproductibles à partir de topologies, de vulnérabilités
  injectées et de postures de sécurité. Il déploie les scénarios, conserve une
  vérité terrain et calcule les mesures d'évaluation.
\end{description}

Une \gls{api} assure le lancement et le suivi des campagnes, tandis qu'une
interface présente la topologie, les événements et les artefacts. La chaîne
d'apprentissage mobilise PyTorch, Transformers, PEFT, bitsandbytes et TRL ; le
déploiement des environnements repose sur des descriptions déclaratives et une
automatisation de l'infrastructure. Les résultats sont enregistrés dans des
formats structurés afin d'être analysés indépendamment de la génération du
modèle.

\section{Limites initiales}

\begin{description}
  \item[Comportement trop généraliste] Un même modèle assumait des tâches aux
  contraintes différentes. Les rôles dépendaient principalement des prompts,
  sans garantir la conformité des appels d'outils ni la stabilité du
  comportement.

  \item[Spécialisation incomplète] Il manquait un processus continu pour
  nettoyer les traces, produire des exemples, adapter le modèle avec des
  ressources limitées, servir plusieurs spécialités et les sélectionner pendant
  une campagne.

  \item[Pipeline insuffisamment contractuel] Une sortie syntaxiquement valide
  pouvait encore porter une hypothèse non vérifiée ou être transmise à la phase
  suivante. Les interruptions et erreurs d'outils devaient aussi être
  représentées sans ambiguïté.

  \item[Mesure partielle] La comparaison à la vérité terrain ne caractérisait
  pas assez finement le comportement de l'agent. Une évaluation sémantique seule
  aurait, à l'inverse, rendu les scores sensibles au modèle juge et à la
  formulation.

  \item[Contexte volumineux] Les topologies, sorties de commandes et observations
  répétées pouvaient saturer la fenêtre du modèle, alourdir le traitement et masquer
  des informations nécessaires à la décision.

  \item[Priorisation topologique limitée] La représentation du réseau décrivait
  les hôtes et leurs relations, mais évaluait peu l'effet d'une nouvelle
  découverte sur les chemins menant à une cible.
\end{description}

\section{Mission principale~: spécialisation \gls{hmoe} agentique}

La mission prioritaire était de construire une architecture \gls{hmoe} adaptée à
LANCE. Contrairement à un mélange d'experts neuronal qui route les jetons, elle
partage un modèle quantifié entre des adaptateurs LoRA spécialisés par rôle ; le
routeur active l'un d'eux selon la phase ou le rôle demandé.

Le travail couvre la préparation, la validation et la diversification des traces,
l'entraînement \gls{qlora}, le service derrière une \gls{api} commune et le
routage. Il comprend aussi l'intégration des experts à l'agent qui planifie,
appelle les outils, interprète leurs réponses et transmet un état structuré, avec
un routage stable, des changements d'adaptateur sérialisés et une compatibilité
avec l'existant.

\section{Pipeline et harness comme protocole de validation}

La spécialisation ne se juge pas sur la seule perte d'entraînement. Le pipeline
sert de protocole opérationnel : il délimite les rôles, injecte les outils,
valide les schémas et les preuves, maintient la cohérence entre phases et signale
toute campagne incomplète.

Le harness fournit le protocole expérimental complémentaire : il crée des
situations contrôlées, associe les faiblesses attendues à des indicateurs et
compare les artefacts à la vérité terrain. Les scénarios, leur déploiement, le
suivi des campagnes et les métriques permettent ainsi de déterminer si la
spécialisation améliore l'action de l'agent au-delà du style de ses réponses.

\section{Objectifs opérationnels}

La mission se décline en trois objectifs~:

\begin{enumerate}
  \item concevoir la préparation des données, l'entraînement \gls{qlora}, le
  service multi-adaptateurs et le routage des experts ;
  \item consolider le pipeline afin de rendre transitions, appels d'outils et
  preuves vérifiables ;
  \item étendre le harness pour déployer des scénarios reproductibles et mesurer
  le comportement des différentes configurations.
\end{enumerate}

Ces objectifs forment un même système : les experts décident, le pipeline encadre
leur exécution et le harness les compare. Les optimisations de contexte et de
graphe soutiennent cet ensemble sans remplacer la spécialisation agentique.

\chapter{Conception et implémentation de la solution}
\label{chap:implementation}

\section{Démarche d'ingénierie}

Mon travail a suivi une démarche itérative fondée sur l'observation des
exécutions plutôt que sur l'ajout isolé de fonctionnalités. Pour chaque
évolution, j'ai d'abord identifié un défaut reproductible dans une trace, un
artefact de phase ou une campagne de benchmark. J'ai ensuite localisé la
frontière responsable --- modèle, contrat d'outil, orchestration ou évaluateur
--- avant d'implémenter un correctif et d'ajouter un test de non-régression.
Cette méthode était particulièrement importante pour évaluer l’agent génératif intégré à LANCE: une sortie conforme au format attendu peut néanmoins être techniquement erronée. Inversement, un échec observé ne provient pas nécessairement du modèle, mais peut résulter d’un contexte d’entrée tronqué ou d’un contrat d’interface ambigu.


Les tests ont porté sur les frontières les plus risquées : validité des données
d'entraînement, compatibilité de l'environnement, sélection de l'expert,
respect des outils exposés, reprise après rejet d'un livrable, cohérence des
artefacts entre phases et calcul des scores. Le but n'était pas seulement de
faire terminer une campagne, mais de rendre explicite la raison pour laquelle
elle pouvait être considérée comme complète.

\section{Architecture générale}

LANCE sépare les décisions probabilistes des contrôles déterministes. La classe
\texttt{\footnotesize Pipeline} orchestre six phases et associe à chacune un prompt, un
livrable attendu et un ensemble d'outils autorisés. Le modèle propose une
analyse ou une action ; l'orchestrateur exécute l'outil, journalise son résultat
et décide si le contrat de la phase est satisfait. Les artefacts structurés
forment ainsi une mémoire interphase plus fiable qu'une conversation libre.

Autour du pipeline, FastAPI expose le lancement et le suivi des campagnes. Un
flux \gls{sse} alimente l'interface, qui affiche la topologie, les livrables et
les métriques. Le benchmark déploie les scénarios au moyen d'Ansible sur
Proxmox, tandis qu'un évaluateur distinct compare les artefacts à la vérité terrain.
Enfin, un serveur \gls{hmoe} compatible avec l'\gls{api} OpenAI fournit au pipeline un point
d'accès unique à plusieurs adaptateurs spécialisés.

La \cref{fig:architecture-lance} présente les principales frontières du système :
le déploiement fournit la cible au pipeline, tandis que la vérité terrain reste
du côté de l'évaluateur. Le pipeline conserve seul la responsabilité d'autoriser
les outils, d'exécuter les actions et de produire les artefacts évalués.

\begin{figure}[htbp]
  \centering
  \resizebox{\textwidth}{!}{\usetikzlibrary{shapes, arrows.meta, positioning, fit}

\begin{tikzpicture}[
  font=\sffamily\small,
  >={Latex[length=3mm, width=2.2mm]},
  flow/.style={->, thick, draw=black!75},
  feedback/.style={->, thick, dashed, draw=violet!80!black},
  component/.style={draw=blue!75!black, fill=blue!8, rounded corners=3pt,
    align=center, minimum height=11mm, text width=30mm, inner sep=5pt, semithick},
  phase/.style={draw=cyan!75!black, fill=cyan!8, rounded corners=2pt,
    align=center, minimum height=10mm, text width=22mm, inner sep=4pt, font=\sffamily\footnotesize, semithick},
  external/.style={draw=orange!80!black, fill=orange!10, rounded corners=3pt,
    align=center, minimum height=11mm, text width=30mm, inner sep=5pt, semithick},
  artifact/.style={draw=green!60!black, fill=green!10, rounded corners=3pt,
    align=center, minimum height=11mm, text width=32mm, inner sep=5pt, semithick},
  group/.style={draw=black!50, dashed, rounded corners=4pt, inner sep=8pt, semithick},
  lbl/.style={font=\sffamily\scriptsize, fill=white, inner sep=1.5pt} % Style pour texte sur les flèches
]

  % --- Briques externes (Gauche) ---
  \node[external] (catalogue) {Catalogue de scénarios\\et vérité terrain};
  \node[external, below=12mm of catalogue] (deploy) {Ansible + Proxmox\\déploiement reproductible};
  \node[external, below=12mm of deploy] (target) {Réseau cible\\IoT/OT};

  % --- Interface & API ---
  \node[component, right=22mm of catalogue] (ui) {Tableau de bord\\opérateur};
  \node[component, below=12mm of ui] (api) {API FastAPI\\suivi par SSE};

  % --- Pipeline LANCE ---
  \node[phase, right=18mm of api, yshift=24mm] (p1) {1. Analyse\\du graphe};
  \node[phase, right=6mm of p1] (p2) {2. Recon-\\naissance};
  \node[phase, right=6mm of p2] (p3) {3. Vulnéra-\\bilités};
  \node[phase, below=10mm of p3] (p4) {4. Vérifi-\\cation};
  \node[phase, left=6mm of p4] (p5) {5. Intrusion\\contrôlée};
  \node[phase, left=6mm of p5] (p6) {6. Rapport\\final};
  
  \node[group, fit=(p1)(p2)(p3)(p4)(p5)(p6),
    label={[font=\sffamily\small\bfseries, yshift=2pt]above:Pipeline LANCE}] (pipeline) {};

  % --- Modèle & Outils ---
  \node[component, right=18mm of p3] (model) {Fournisseur LLM\\HMoE local ou API distante};
  \node[component, below=12mm of model] (tools) {Registre d'outils\\Nmap, HTTP, MQTT,\\Modbus, SSH, \ldots};

  % --- Artefacts & Évaluation ---
  \node[artifact, below=18mm of p6, xshift=15mm] (artifacts) {Artefacts structurés\\et traces d'appels};
  \node[artifact, right=18mm of artifacts] (harness) {Harness strict-v3\\comparaison à la vérité terrain};
  \node[artifact, right=18mm of harness] (metrics) {Scores, diagnostics\\et rapports de campagne};

  % --- Liens et Flux ---
  \draw[flow] (catalogue) -- (deploy);
  \draw[flow] (deploy) -- (target);
  \draw[flow] (ui) -- (api);
  \draw[flow] (api.east) -- (pipeline.west);
  
  \draw[flow] (p1) -- (p2);
  \draw[flow] (p2) -- (p3);
  \draw[flow] (p3) -- (p4);
  \draw[flow] (p4) -- (p5);
  \draw[flow] (p5) -- (p6);
  
  \draw[flow, bend left=15] (pipeline.east) to node[lbl, above=1pt]{requête} (model.west);
  \draw[feedback, bend left=15] (model.west) to node[lbl, below=1pt]{réponse / appel} (pipeline.east);
  
  \draw[flow] (pipeline.south east) -- (tools.north west);
  \draw[flow] (tools.west) -- node[lbl, below]{actions bornées} (target.east);
  \draw[feedback] (target.north east) -- node[lbl, above]{observations} (pipeline.south west);
  
  \draw[flow] (pipeline.south) -- (artifacts.north);
  \draw[flow] (artifacts) -- (harness);
  \draw[flow] (harness) -- (metrics);
  
  \draw[flow] (catalogue.south east) |- node[lbl, pos=0.8, above]{référence scellée} (harness.north);
  \draw[feedback] (metrics.north west) |- (ui.south);

\end{tikzpicture}
}
  \caption{Architecture fonctionnelle de LANCE et de son environnement d'évaluation}
  \label{fig:architecture-lance}
\end{figure}


\section{Mission principale : construction d'un HMoE agentique}

\subsection{Principe de spécialisation}

La mission principale de mon stage a consisté à construire une chaîne HMoE
agentique, depuis les traces d'exécution jusqu'au déploiement des experts dans
LANCE. Ici, le terme HMoE ne désigne pas une couche de mélange d'experts apprise
à l'intérieur du réseau. Le système partage un modèle de base, charge au
démarrage quatre adaptateurs \gls{peft} spécialisés, puis active celui qui
correspond à chaque requête : \texttt{\footnotesize secretary}, \texttt{\footnotesize recon},
\texttt{\footnotesize vuln} et \texttt{\footnotesize exploit}. Cette organisation reprend
l'idée de spécialisation des mélanges d'experts tout en l'adaptant aux
contraintes matérielles du projet \cite{shazeer2017moe}.

Le secrétaire traite l'analyse initiale du graphe et la synthèse finale. L'expert
de reconnaissance couvre la phase~2 ; l'expert de vulnérabilités intervient en
phase~3 ; l'expert d'exploitation couvre les phases~4 et~5. Cette répartition
évite d'entraîner chaque adaptateur sur l'intégralité des outils et des formats.
Elle réduit aussi l'ambiguïté : une trace d'exploitation et une trace de rapport
ne demandent ni les mêmes actions ni le même critère de terminaison.

Le choix d'adaptateurs plutôt que de modèles indépendants répond à une
contrainte de mémoire GPU. Une seule instance quantifiée du modèle de base est
chargée, puis l'adaptateur correspondant est sélectionné pour chaque requête. Le gain de
mémoire a une contrepartie : le changement d'adaptateur modifie un état partagé.
J'ai donc sérialisé la génération avec un verrou. Les requêtes peuvent être
reçues concurremment par l'API, mais deux générations ne changent jamais
l'adaptateur au même instant.

La \cref{fig:architecture-hmoe} détaille ce routage. La sélection porte sur une
requête complète et précède la génération : les quatre experts partagent donc
les poids quantifiés du même modèle, mais appliquent chacun une correction
QLoRA propre à leur rôle.

\begin{figure}[htbp]
  \centering
  \resizebox{0.98\textwidth}{!}{\begin{tikzpicture}[
  font=\sffamily\small,
  >={Latex[length=2.8mm,width=2mm]},
  flow/.style={->,thick,draw=black!75},
  link/.style={thick,draw=black!65},
  return/.style={->,thick,dashed,draw=violet!80!black},
  block/.style={draw=blue!70!black,fill=blue!8,rounded corners=3pt,
    align=center,minimum height=11mm,text width=31mm,inner sep=5pt,semithick},
  expert/.style={draw=green!55!black,fill=green!9,rounded corners=3pt,
    align=center,minimum height=10mm,text width=24mm,inner sep=4pt,semithick},
  guard/.style={draw=orange!80!black,fill=orange!9,rounded corners=3pt,
    align=center,minimum height=11mm,text width=94mm,inner sep=5pt,semithick},
  group/.style={draw=black!50,dashed,rounded corners=4pt,inner sep=7pt,semithick},
  lbl/.style={font=\sffamily\scriptsize,fill=white,inner sep=1.5pt}
]
  % Une requête entière est routée avant toute génération
  \node[block] (pipeline) {Pipeline LANCE\\requête OpenAI-compatible};
  \node[block,right=20mm of pipeline] (router)
    {Routeur déterministe\\phase et rôle\\jamais par jeton};
  \coordinate (routingmid) at ($(pipeline)!0.5!(router)$);
  \coordinate (split) at ($(routingmid)+(0,-18mm)$);

  % Quatre adaptateurs spécialisés clairement séparés
  \node[expert,below=20mm of routingmid,xshift=-54mm] (secretary)
    {\texttt{secretary}\\phases 1 et 6};
  \node[expert,right=8mm of secretary] (recon)
    {\texttt{recon}\\phase 2};
  \node[expert,right=8mm of recon] (vuln)
    {\texttt{vuln}\\phase 3};
  \node[expert,right=8mm of vuln] (exploit)
    {\texttt{exploit}\\phases 4 et 5};
  \node[group,fit=(secretary)(recon)(vuln)(exploit),
    label={[font=\sffamily\small\bfseries]above:Adaptateurs QLoRA spécialisés}] (adapters) {};

  % Socle commun et exécution protégée
  \node[block,below=18mm of adapters] (base)
    {Qwen2.5-3B-Instruct\\base quantifiée 4 bits\\partagée et gelée};
  \node[block,below=14mm of base] (generation)
    {Adaptateur actif\\+ génération sérialisée};
  \node[guard,below=14mm of generation] (guards)
    {Garde-fous : budget de contexte et de génération, filtrage des outils\\
     autorisés et reprise contrôlée après un rejet};

  % Routage de la requête
  \draw[flow] (pipeline) -- node[above,font=\scriptsize]{requête complète} (router);
  \draw[flow] (router.south) |- (split);
  \draw[flow] (split) -| (secretary.north);
  \draw[flow] (split) -| (recon.north);
  \draw[flow] (split) -| (vuln.north);
  \draw[flow] (split) -| (exploit.north);

  % Convergence vers l'unique modèle de base
  \coordinate (merge) at ($(base.north)+(0,8mm)$);
  \draw[link] (secretary.south) |- (merge);
  \draw[link] (recon.south) |- (merge);
  \draw[link] (vuln.south) |- (merge);
  \draw[link] (exploit.south) |- (merge);
  \draw[flow] (merge) -- (base.north);
  \draw[flow] (base) -- (generation);
  \draw[return] (guards.north) -- node[lbl,right]{contraintes} (generation.south);
  \draw[return] (generation.west) -- ++(-62mm,0) |-
    node[lbl,pos=0.25,below]{texte ou appel d'outil} (pipeline.south);
\end{tikzpicture}
}
  \caption{Architecture du HMoE agentique et routage déterministe par phase}
  \label{fig:architecture-hmoe}
\end{figure}

\subsection{Construction des données d'apprentissage}

La qualité du système dépend d'abord de la correspondance entre les données et
la boucle agentique réellement exécutée. J'ai transformé les traces en exemples
de conversation contenant les messages, les appels d'outils, leurs résultats et
les schémas d'outils disponibles. Les exemples sont routés vers l'expert associé
à la phase qui a produit le livrable. Cette conservation des appels est
essentielle : un corpus qui ne garderait que la réponse finale apprendrait à
rédiger, mais pas à décider quand observer le réseau ou quand sauvegarder un
artefact.

La préparation cible Qwen2.5-3B-Instruct et applique le gabarit de conversation
utilisé à l'inférence. Pour chaque exemple, seuls les outils effectivement
appelés et quelques distracteurs sont conservés. Les distracteurs obligent le
modèle à sélectionner un outil pertinent sans gonfler inutilement le contexte.
Les conversations longues sont découpées à la frontière des tours assistant ;
les couples appel/résultat restent liés. La phase~5 constitue une exception : sa
chaîne de preuves est gardée atomique afin qu'une conclusion d'intrusion ne soit
jamais séparée de l'action qui la soutient.

Lorsque la fenêtre reste trop longue, la préparation réduit d'abord les sorties
d'outils volumineuses, puis les arguments textuels et enfin le milieu des
messages système. Le préfixe contenant le rôle et le suffixe décrivant le
critère de terminaison sont préservés. Chaque ligne indique dans ses métadonnées
le modèle cible, la fenêtre source et l'éventuelle troncature. Les limites sont
alignées sur l'usage futur : 6\,144 tokens pour le secrétaire, la reconnaissance
et l'exploitation, et 4\,096 pour l'analyse de vulnérabilités.

J'ai également intégré une boucle de correction supervisée. Le harness peut
extraire des candidats correspondant à un faux négatif, un faux positif, une
sévérité incorrecte ou un livrable incohérent. Ces candidats restent en attente
jusqu'à une revue humaine ; seuls ceux explicitement acceptés sont convertis en
exemples SFT. Ils sont attribués à la première phase disposant d'un signal causal
observable. Une cible absente de la reconnaissance corrige ainsi l'expert
\texttt{\footnotesize recon}, et non automatiquement l'expert \texttt{\footnotesize vuln}. Ce choix évite
d'enseigner à un composant aval à compenser silencieusement une erreur amont.

\subsection{Préflight et reproductibilité}

Avant tout entraînement, un préflight vérifie les conditions sans charger les
poids ni lancer d'optimisation. Il contrôle la structure YAML, la présence du
gabarit et de ses masques de génération, puis parcourt les JSONL. Chaque ligne
doit contenir des messages, au moins une réponse assistant et des métadonnées
cohérentes avec l'expert et le modèle préparé. Toute référence au split
\texttt{\footnotesize eval-sealed} est refusée afin d'empêcher une contamination des données
d'évaluation.

Le préflight contrôle aussi les versions de PyTorch, Transformers, TRL, PEFT,
bitsandbytes, Datasets et Accelerate. Il vérifie que la signature de
\texttt{\footnotesize SFTConfig} expose les paramètres requis par le script, que CUDA est
disponible, que le GPU accepte BF16 et, en mode strict, qu'aucun autre processus
de calcul ne l'occupe. Il produit un rapport JSON comprenant les tailles et
empreintes des jeux de données. Lors d'une promotion, le même outil exige pour
chaque expert un fichier de configuration, les poids finaux et le tokenizer,
puis confirme que l'adaptateur référence le modèle de base prévu.


\subsection{Entraînement QLoRA}

J'ai retenu QLoRA pour adapter le modèle avec une empreinte mémoire compatible
avec la machine disponible \cite{dettmers2023qlora}. Qwen2.5-3B-Instruct est
chargé en 4 bits avec la quantification NF4, une double quantification et des
calculs en BF16. Les adaptateurs ont un rang \(r=16\), un facteur
\(\alpha=32\) et un dropout de 0,05. Ils ciblent les projections
\texttt{\footnotesize q\_proj}, \texttt{\footnotesize k\_proj}, \texttt{\footnotesize v\_proj} et
\texttt{\footnotesize o\_proj}, ainsi que \texttt{\footnotesize gate\_proj}, \texttt{\footnotesize up\_proj} et
\texttt{\footnotesize down\_proj} dans le bloc feed-forward.

L'entraînement utilise un lot unitaire et une accumulation de gradient pour
obtenir un lot effectif plus grand sans dépasser la mémoire disponible. Le
gradient checkpointing, l'optimiseur paginé et la désactivation du cache
complètent ce dispositif. Le jeu est séparé de manière déterministe entre
apprentissage et validation. La perte n'est calculée que sur les tokens produits
par l'assistant grâce aux masques du gabarit ; les messages système, demandes et
résultats d'outils servent de contexte, mais ne sont pas des cibles à recopier.


\subsection{Routage et état d'exécution sans session}

Le serveur expose cinq identifiants : un identifiant HMoE automatique et un
identifiant de diagnostic par expert. Le routage automatique inspecte la phase
et le rôle contenus dans la requête système. Le routage explicite reste utile
pour tester un adaptateur isolément, mais le tableau de bord ne doit connaître
que l'identifiant automatique. La spécialisation demeure ainsi interne au
service et ne complexifie pas la configuration d'une campagne.

Le modèle partage un état de poids, mais le suivi agentique est volontairement
sans session. À chaque requête, le serveur reconstruit l'avancement à partir de
la conversation OpenAI reçue : nombre d'appels réussis, tentatives de sauvegarde
rejetées, exigences encore manquantes et progression de l'intrusion. Il injecte
ensuite ce résumé autoritaire dans le dernier message. Aucune mémoire de campagne
n'est conservée entre deux requêtes, ce qui empêche le mélange de deux exécutions
concurrentes et facilite la reprise après redémarrage du service.

\subsection{Contrats d'outils, compactage et récupération}

La petite taille des experts rend les boucles d'outils sensibles au bruit. Pour
la reconnaissance, seuls cinq schémas sont rendus dans le prompt : découverte
ARP, découverte Nmap, scan Nmap, lecture et sauvegarde de livrable. Le registre
exécutable complet reste chez l'appelant ; la réduction concerne seulement ce
que voit le modèle. Dès que la couverture minimale est atteinte, le schéma visible
est réduit à la sauvegarde. À l'inverse, tout appel généré vers un outil absent
du contrat de la requête est supprimé avant d'être transmis au pipeline.

Le compactage est aligné sur les fenêtres d'entraînement. Le serveur remplace
les brouillons de livrables rejetés, résume les anciennes sorties réseau et, si
nécessaire, replie les vieux tours en un registre déterministe. Il préserve les
paires appel/résultat récentes et la dernière erreur contractuelle. La génération
est ensuite plafonnée selon l'expert et l'espace restant dans la fenêtre. Ce
mécanisme ne cherche donc pas seulement à réduire les tokens : il conserve les
informations dont dépend la prochaine action.

Lorsqu'une tentative de sauvegarde retourne une liste structurée
\texttt{\footnotesize missing\_requirements}, le serveur peut produire directement l'appel
manquant autorisé. Il transforme, par exemple, un port non couvert en scan Nmap
ciblé. Cette récupération déterministe intervient avant une nouvelle génération
et évite que le modèle répète le même livrable refusé. Dans les phases
d'intrusion, le pipeline maintient en plus un registre d'actions faisant autorité.
Si le modèle s'arrête ou si le fournisseur local devient indisponible, une
reprise bornée exécute uniquement des sondes déduites du contexte, sans inventer
d'identifiants. Le livrable final est alors reconstruit à partir du registre et
signale les couvertures encore incomplètes.

\subsection{Déploiement du service}

Le service HMoE charge le modèle de base quantifié et les quatre adaptateurs,
puis expose une API compatible OpenAI avec des routes de santé, de modèles et de
conversation. Il est lancé comme service utilisateur systemd. Sur la machine
GPU, l'écoute est limitée au pont Docker pour OpenWebUI ; l'accès depuis
\texttt{\footnotesize nato-master} passe par un proxy lié uniquement à l'adresse Tailscale.
Cette topologie évite d'exposer publiquement une API qui ne valide pas elle-même
un jeton Bearer. Après une période d'inactivité configurable, les allocations
CUDA inutilisées sont rendues au pilote sans interrompre le modèle chargé.


\section{Évolution détaillée du pipeline}

\subsection{Phase 2 : reconnaissance contractuelle}

La phase~2 transforme une cible réseau en surface d'attaque observée. L'agent
doit couvrir la découverte, les hôtes attendus et leurs services, puis enregistrer
\texttt{\footnotesize 02\_recon.md}. La sauvegarde n'est pas un simple signal de fin : l'outil
vérifie les appels effectués et renvoie les exigences manquantes. Les résultats
de scan sont projetés dans une représentation déterministe des hôtes, ports et
protocoles. Cette projection alimente la phase suivante même si la prose du
livrable est incomplète.

J'ai ajouté une réconciliation entre cette surface et le graphe de scénario.
Elle associe les observations aux nœuds connus et conserve les cibles non
résolues au lieu de les faire disparaître. En mode découverte, les liens peuvent
être inférés après la reconnaissance. Ainsi, la phase~3 ne dépend plus uniquement
d'une topologie préexistante ni de la capacité du modèle à recopier exactement
les résultats Nmap.

\subsection{Phase 3 : analyse par appareil}

La phase~3 est décomposée en trois étapes. Un scanner déterministe collecte
d'abord les indices disponibles pour chaque appareil. Des micro-agents spécialisés
analysent ensuite cette surface appareil par appareil, avec une limite de temps,
de tours et de tokens. Enfin, une agrégation déterministe fusionne les sorties
dans une file canonique de vulnérabilités. Une défaillance locale ne supprime pas
les résultats du scanner : elle est inscrite dans un artefact de statut et la
campagne continue avec un résultat partiel explicite.

L'analyse conserve les versions de logiciels et peut interroger un catalogue CVE
hors ligne. Les findings sont normalisés avec un identifiant stable, une cible,
un type, une sévérité, une preuve et un plan de vérification. Dans le profil HMoE
compact, certaines observations directes peuvent être retenues pour le rapport
sans être promues automatiquement vers l'exploitation. Ce cloisonnement évite
qu'une configuration seulement suspectée devienne une action offensive.

\subsection{Phase 4 : vérification par vulnérabilité}

La phase~4 planifie un travail indépendant pour chaque finding éligible. Le
contrôleur choisit un outil compatible avec le service et impose un contrat de
vérification. Le micro-agent ne reçoit donc pas une liberté d'exploitation
illimitée : il doit exécuter une sonde ciblée, rester dans le périmètre autorisé
et produire un résultat lié à cette sonde. Dans le profil compact, si l'appel au
modèle échoue ou si la sonde requise n'apparaît pas, le contrôleur peut exécuter
une sonde de repli bornée.

L'agrégateur distingue un succès, un échec conclusif, une erreur et l'absence de
test. Un succès enrichit la preuve ; un échec peut réfuter une simple suspicion ;
une erreur ne transforme pas la détection en faux positif. Les résultats sont
réunis dans \texttt{\footnotesize 04\_exploitation.json} avec leur provenance d'outil. Les
hôtes découverts durant cette phase déclenchent un cycle limité de reconnaissance
et d'analyse avant la phase~5, ce qui permet de traiter un pivot sans relancer
toute la campagne.

\subsection{Phase 5 : intrusion contrôlée}

La phase~5 raisonne sur les findings vérifiés, les chemins possibles et les
identifiants réellement récupérés. Un contexte structuré énumère les points
d'entrée, les cibles, les services et les chaînes candidates. Le contrat terminal
exige des actions observables : lecture du contexte, couverture des points
d'entrée, essais de réutilisation des identifiants sur les cibles concernées et,
lorsqu'un identifiant est disponible, au moins un résultat d'accès authentifié.

Chaque action est inscrite dans un registre faisant autorité. Le modèle produit
un mémo, mais il ne peut pas déclarer seul la phase complète. Le contrôleur
calcule la couverture, propose une action lorsqu'aucun progrès n'est observé et
limite les reprises à un nombre borné de tours. Après un accès authentifié, une
collecte post-accès récupère uniquement les éléments nécessaires à la preuve.
Une campagne qui n'exécute aucune action ou ne satisfait pas le contrat est
marquée bloquée ou incomplète ; elle n'est pas silencieusement convertie en
succès.

\subsection{Phase 6 : rapport fondé sur les artefacts}

La phase~6 ne réinjecte pas l'intégralité des traces. Le pipeline construit un
contexte compact à partir des vulnérabilités, résultats de vérification, chaînes
d'intrusion et observations de configuration. Les tables lourdes sont
pré-générées, et seuls les findings confirmés avec un niveau de preuve suffisant
sont présentés comme vérifiés. Le modèle se concentre ainsi sur l'analyse et la
hiérarchisation plutôt que sur la recopie de données volumineuses.

Le rapport généré est validé avant fusion avec les sections déterministes. Si le
modèle local échoue ou produit une synthèse incompatible avec les artefacts, un
repli compose le rapport à partir des données structurées. Cette dernière
frontière garantit qu'une formulation persuasive ne réintroduit pas une
vulnérabilité éliminée par la phase~4.

\section{Harness expérimental et politique d'évaluation \texttt{\footnotesize strict-v3}}

Le harness expérimental organise le déploiement, la collecte des artefacts et
la répétition des campagnes. Au sein de ce dispositif, l'évaluateur applique
la politique \texttt{\footnotesize strict-v3} pour mesurer non seulement la détection, mais aussi la
qualité de la correspondance et la traçabilité des preuves. Chaque scénario
associe une topologie, des packs de vulnérabilités et une vérité terrain. Des
contrats de matching revus décrivent les champs structurants et les
compatibilités autorisées. Le moteur construit les correspondances candidates,
puis résout une association bipartite globale : un finding ne peut expliquer
qu'une entrée de vérité terrain, et réciproquement. Cette règle supprime les
scores artificiels où une même observation créditerait plusieurs vulnérabilités.

Dans cette nomenclature, \texttt{\footnotesize strict-v3} désigne la politique de scoring,
\texttt{\footnotesize strict-v3.3} le contrat métrique, \texttt{\footnotesize strict-v3.2} le schéma des
contrats de matching et \texttt{\footnotesize evidence-v2} le contrat de preuve. Ces
versions séparées permettent de faire évoluer un contrat sans changer
implicitement le sens des autres résultats.

Le crédit dépend de la précision structurelle du match. Une correspondance
exacte reçoit le crédit maximal ; un type correct mais incomplet ou une
compatibilité explicitement prévue reçoit un crédit moindre. Le score ajusté à
la qualité combine ce crédit avec l'erreur de sévérité et la qualité de
vérification. Le \emph{verified F1} impose en outre un appel d'outil lié au
finding et dont le résultat est explicitement positif. Les métriques de chemins
exigent que toutes les vulnérabilités attendues soient retrouvées ; leurs
variantes vérifiées contrôlent aussi l'ordre des appareils dans une chaîne de
phase~5.

J'ai appliqué une politique d'échec fermé aux métadonnées et aux artefacts. Une
campagne incomplète reste incomplète. Une version de contrat métrique incompatible
neutralise les scores dépendant des preuves au lieu de les calculer avec une
sémantique différente. Pour le split scellé, le worker ne reçoit ni vérité
terrain ni détails privés ; le contrôleur renvoie seulement un résumé signé. Les
scénarios attendus mais absents valent zéro lors de l'agrégation officielle.

L'agrégation évite enfin qu'un scénario riche en findings domine le résultat.
Les répétitions sont d'abord moyennées pour chaque scénario, puis les scénarios
sont macro-moyennés à poids égal. Les métriques brutes, créditées,
ajustées et vérifiées restent séparées afin qu'une amélioration de précision ne
masque pas une régression de preuve. Un juge sémantique peut aider au diagnostic,
mais la correspondance un-à-un, la validation de schéma et les scores finaux
restent déterministes, ce qui limite les biais du paradigme
\emph{LLM-as-a-judge} \cite{zheng2023judge}.

\section{Maîtrise du contexte}

J'ai borné les sorties SSH, compacté les résumés de topologie et dédoublonné
certaines observations MQTT ou HTTP. Ces optimisations ont été appliquées avec
prudence : une première réduction trop agressive de la surface d'attaque a été
annulée parce qu'elle supprimait une information utile. Ce retour a directement
guidé le compactage du serveur HMoE, qui conserve désormais les contrats, les
erreurs récentes et le registre des outils avant de réduire la prose.

\paragraph{Exploration topologique initiale.}
Au début du stage, j'ai également relié au pipeline un graphe d'attaque orienté
pondéré. Le coût d'une arête est l'inverse d'un poids heuristique de
franchissement. Après une découverte ou une compromission, les distances vers la
cible sont recalculées ; leur variation localise les nœuds rapprochés de
l'objectif. Une propagation sur les prédécesseurs forme un cône d'impact, dans
l'esprit des graphes d'attaque multi-hôtes \cite{ou2005mulval}. Ces valeurs ne
sont pas des probabilités calibrées : elles servent uniquement à prioriser une
analyse. Cette expérimentation est restée un apport secondaire et ne constitue
pas l'axe principal des campagnes présentées dans ce rapport.

\section{Résultats et difficultés}
\subsection{Instantané expérimental du 24 août 2026}
\label{sec:resultats-provisoires}

Les résultats présentés ici constituent un instantané de travail et non une
campagne finale. J'ai retenu quatre exécutions terminées et comparables au
regard des métadonnées disponibles, accessibles sur \texttt{nato-master}.
Elles utilisent toutes le modèle MiniMax-M2.7, le profil \texttt{full}, le
split \texttt{dev-public}, la révision logicielle \texttt{e7c4f6a} et la
politique d'évaluation \texttt{strict-v3} avec le contrat métrique
\texttt{strict-v3.3}. Les métadonnées indiquent que le modèle n'avait pas accès
à la vérité terrain. Chaque scénario ne comporte qu'une répétition ; les
valeurs ne permettent donc pas encore d'estimer la variance.

\begin{table}[H]
  \centering
  \small
  \caption{Résultats provisoires des quatre campagnes sélectionnées}
  \label{tab:resultats-provisoires}
  \begin{tabular}{lrrrrr}
    \toprule
    Scénario & Précision & Rappel & F1 & F1 vérifié & F1 qualité \\
    \midrule
    S1 & \(0{,}875\) & \(0{,}583\) & \(0{,}700\) & \(0{,}600\) & \(0{,}625\) \\
    S2 & \(0{,}833\) & \(0{,}769\) & \(0{,}800\) & \(0{,}640\) & \(0{,}670\) \\
    S5 & \(0{,}900\) & \(0{,}600\) & \(0{,}720\) & \(0{,}640\) & \(0{,}680\) \\
    S8 & \(0{,}533\) & \(0{,}571\) & \(0{,}552\) & \(0{,}414\) & \(0{,}448\) \\
    \midrule
    Moyenne macro & \(0{,}785\) & \(0{,}631\) & \(0{,}693\) & \(0{,}574\) & \(0{,}606\) \\
    \bottomrule
  \end{tabular}
\end{table}

Comme l'indique \cref{tab:resultats-provisoires}, la moyenne macro atteint un
F1 de \(0{,}693\), contre \(0{,}574\) lorsque seuls les constats vérifiés sont
crédités. Le F1 ajusté par la qualité s'établit à \(0{,}606\). L'écart entre
ces mesures suggère qu'une détection correcte au niveau du type et de la cible
ne s'accompagne pas toujours d'une preuve d'exécution suffisamment forte.

La complétion moyenne de la phase 4 atteint \(92{,}8\,\%\) et la couverture
d'exploitation \(82{,}4\,\%\). En revanche, la couverture macro des chemins
n'est que de \(52{,}1\,\%\), et celle des chemins entièrement vérifiés de
\(6{,}3\,\%\). S1 est le seul scénario de cet échantillon à obtenir une
couverture vérifiée non nulle, égale à \(25\,\%\). Dans cet échantillon, la
difficulté principale semble donc être la conservation d'une chaîne de preuves
complète sur plusieurs étapes, et non le seul lancement des vérifications
locales.

S2 obtient le meilleur F1 de l'échantillon avec \(0{,}800\). S8 est le cas le
plus difficile : huit vrais positifs, sept faux positifs et six faux négatifs
conduisent à un F1 de \(0{,}552\). Ce résultat motive les gardes consacrées à
la précision sur les scénarios étendus et à la réduction des hallucinations.

Les quatre exécutions totalisent environ 3,58 dollars, 11,0 millions de tokens,
1\,202 appels d'outils et 30 erreurs d'outils. Elles utilisent cependant un
modèle généraliste distant et non le \gls{hmoe} local. Cet instantané constitue
donc un test d'intégration du pipeline et du harness ; il ne démontre
pas encore l'efficacité des experts entraînés. Une campagne finale devra figer
la version du code, répéter chaque scénario et publier moyenne, dispersion et
intervalle de confiance pour les baselines et le \gls{hmoe}.

Les identifiants de runs utilisés sont
\texttt{2026-08-24\_092203} pour S1,
\texttt{2026-08-21\_144141} pour S2,
\texttt{2026-08-24\_083553} pour S5 et
\texttt{2026-08-21\_151305} pour S8.


\subsection{Difficultés techniques rencontrées}

\begin{description}
  \item[Aligner apprentissage et exécution] Un adaptateur entraîné sur un outil,
  un format ou une fenêtre différents de la production peut échouer malgré une
  perte d'entraînement correcte. J'ai donc partagé le gabarit, les schémas et les
  budgets entre préparation, préflight et service.

  \item[Préserver la preuve sous contrainte de contexte] Une réduction aveugle
  économise des tokens mais peut séparer une conclusion de son appel d'outil. Le
  découpage atomique de la phase~5 et le registre déterministe répondent à ce
  risque.


  \item[Distinguer terminaison et succès] Un livrable sauvegardé ou une réponse
  fluide ne prouve pas la couverture. Les contrats de phase, les reprises bornées
  et l'échec fermé rendent les lacunes visibles au lieu de les masquer.

\end{description}

\chapter{Spécialisation QLoRA et HMoE agentique}
\label{chap:science}

\section{Une mission centrée sur la spécialisation de l'agent}
\label{sec:hmoe-probleme}

La mission scientifique principale du stage a consisté à étudier une
spécialisation frugale de LANCE. Un \gls{llm} généraliste peut expliquer une
vulnérabilité, mais l'agent attendu doit enchaîner des actions beaucoup plus
contraintes : choisir un outil autorisé, produire des arguments conformes à son
schéma, exploiter les observations obtenues et enregistrer le livrable propre à
la phase courante. Ces comportements ne sollicitent pas exactement les mêmes
compétences lors de la reconnaissance, de l'analyse de vulnérabilités, de
l'exploitation ou de la synthèse. Les inscrire uniquement dans le prompt
allonge le contexte et ne garantit pas leur application ; entraîner quatre
modèles complets multiplierait au contraire les besoins en mémoire et en
stockage.

L'hypothèse étudiée est donc la suivante : 

\begin{displayquote}
  un modèle de base quantifié, partagé par plusieurs adaptateurs spécialisés et
  associé à un routage explicite par phase, peut améliorer la conformité
  opérationnelle de l'agent à coût d'entraînement contenu.
\end{displayquote}

Cette hypothèse porte sur le système agentique complet. Elle ne se réduit ni à
la convergence de l'entraînement, ni à la capacité du modèle à compléter une
trace déjà connue. Elle suppose que l'expert sélectionné généralise à des
topologies nouvelles et conserve la capacité d'appeler les outils au bon
moment. Le travail a ainsi couvert toute la chaîne : constitution des corpus,
adaptation \gls{qlora}, service multi-adaptateurs, routage, garde-fous de
contexte et préparation d'un protocole d'évaluation.

\section{Adaptation frugale par QLoRA}
\label{sec:qlora-methode}

La méthode \gls{qlora} conserve les poids du modèle de base gelés et les charge
en quatre bits. Les gradients traversent cette représentation, mais seuls des
adaptateurs de faible rang sont mis à jour. NF4, la double quantification et les
optimiseurs paginés réduisent notamment l'empreinte mémoire de l'apprentissage
\cite{dettmers2023qlora}. Pour une projection de poids \(W\), l'expert \(e\)
apprend la correction

\[
  h = Q_{4}(W)x + \frac{\alpha}{r} B_e A_e x,
\]

où \(Q_{4}(W)\) désigne le poids de base quantifié, \(A_e\) et \(B_e\) les deux
matrices entraînables et \(r\) leur rang. La configuration retenue utilise
\(r=16\), \(\alpha=32\) et un dropout de \(0{,}05\). Les corrections sont
appliquées aux projections de l'attention ainsi qu'aux projections montante,
descendante et de porte du bloc feed-forward. Le calcul s'effectue en
\texttt{bfloat16}, avec accumulation de gradient et activation du
checkpointing de gradient afin de rester compatible avec une carte graphique
de 16\,Gio.

Chaque exemple du corpus de l'expert \(e\) est une conversation
\((x^{(i)},y^{(i)})\in\mathcal{D}_e\), comprenant le contexte, les outils exposés
et la réponse attendue. Un masque \(m_t\) exclut les messages utilisateur et les
résultats d'outils : l'objectif porte uniquement sur les jetons produits par
l'assistant. La fonction minimisée est ainsi

\[
  \mathcal{L}_e =
  -\frac{1}{\sum_{i,t}m_t^{(i)}}
  \sum_{i,t}m_t^{(i)}
  \log p_{\theta,A_e,B_e}
  \left(y_t^{(i)}\mid x^{(i)},y_{<t}^{(i)}\right).
\]

Cette formulation évite d'apprendre à prédire les entrées du protocole et
concentre l'adaptation sur les décisions et productions de l'agent. Une
séparation déterministe conserve 2\,\% de chaque corpus pour l'évaluation. Les
longueurs maximales sont bornées à 6\,144 jetons, sauf pour l'expert de
vulnérabilités limité à 4\,096 jetons. Le prétraitement découpe les traces trop
longues et conserve les outils pertinents, car entraîner sur un contexte plus
large que celui servi créerait un décalage entre apprentissage et inférence.

\section{HMoE agentique : un routage de rôles, pas de jetons}
\label{sec:hmoe-routage}

Le terme \gls{hmoe} désigne ici une architecture agentique hétérogène. Il faut
la distinguer d'un mélange d'experts neuronal classique. Dans ce dernier, une
porte apprise sélectionne, pour chaque jeton et à l'intérieur d'une couche, un
sous-ensemble de blocs feed-forward ; le calcul conditionnel augmente alors la
capacité du réseau sans activer tous ses paramètres \cite{shazeer2017moe}. Dans
LANCE, le modèle neuronal reste unique : le routage intervient au niveau d'une
requête complète et choisit un adaptateur \gls{peft}. Il ne produit donc ni
pondération douce entre experts, ni décision différente pour chaque jeton.

Le routeur exploite d'abord l'identifiant de phase inscrit dans le message
système. Les phases~1 et~6 sont confiées au secrétaire, la phase~2 à la
reconnaissance, la phase~3 à l'analyse de vulnérabilités, et les phases~4 et~5 à
l'exploitation. En l'absence de numéro explicite, des marqueurs de rôle
stabilisent la sélection ; si aucun motif n'est reconnu, le modèle de base sert
de repli. Les identifiants directs des experts restent exposés pour le
diagnostic, tandis que l'identifiant commun \texttt{lance-moe} réalise le
routage automatique. La \cref{tab:hmoe-experts} résume cette décomposition.

\begin{table}[H]
  \centering
  \caption{Décomposition fonctionnelle du \gls{hmoe} agentique}
  \label{tab:hmoe-experts}
  \begin{tabular}{lp{0.16\textwidth}p{0.48\textwidth}}
    \toprule
    Expert & Phases & Compétence principalement apprise \\
    \midrule
    Secrétaire & 1 et 6 & Interprétation de la topologie, continuité des livrables et synthèse finale \\
    Reconnaissance & 2 & Planification des scans, appels d'outils et consolidation des services observés \\
    Vulnérabilités & 3 & Association entre preuves, services et vulnérabilités sans inventer de constat \\
    Exploitation & 4 et 5 & Vérification contrôlée, conservation des preuves et progression vers l'intrusion \\
    \bottomrule
  \end{tabular}
\end{table}

Cette architecture rend la décision de routage interprétable et reproductible.
Elle permet aussi de remplacer un expert sans réentraîner les autres. En
contrepartie, la frontière entre phases est imposée : une requête ambiguë ne
combine pas les compétences de deux adaptateurs. Le changement d'adaptateur
actif est en outre protégé par un verrou de génération ; plusieurs exécutions
concurrentes sont sérialisées afin d'éviter qu'une requête poursuive sa
génération avec les poids d'un autre expert. Des budgets de contexte et de
génération propres aux rôles bornent enfin le coût, mais peuvent tronquer une
trace utile. Ces compromis relèvent de l'architecture agentique et ne sont pas
visibles dans la seule perte d'apprentissage.

\section{Données et observations provisoires}
\label{sec:hmoe-donnees}

Les quatre corpus préparés proviennent de traces structurées selon le protocole
de LANCE. Ils ont été séparés par rôle afin qu'un adaptateur n'apprenne pas des
contrats de sortie incompatibles. Le prévol vérifie le modèle préparé, le rôle
déclaré, les limites de contexte et l'intégrité des fichiers. Les volumes et les
pertes de validation disponibles au moment de la rédaction figurent dans la
\cref{tab:hmoe-apprentissage}.

\begin{table}[H]
  \centering
  \caption{État provisoire des corpus et de l'apprentissage QLoRA}
  \label{tab:hmoe-apprentissage}
  \begin{tabular}{lrrr}
    \toprule
    Expert & Exemples & Contexte maximal & Perte d'évaluation \\
    \midrule
    Secrétaire & 19\,065 & 6\,144 & 0,0421 \\
    Reconnaissance & 9\,257 & 6\,144 & 0,2102 \\
    Vulnérabilités & 6\,951 & 4\,096 & 0,00475 \\
    Exploitation & 14\,629 & 6\,144 & 0,00836 \\
    \bottomrule
  \end{tabular}
\end{table}

Ces valeurs montrent seulement que la vraisemblance des réponses attendues est
mesurable sur la partition réservée. Elles ne sont pas directement comparables
entre experts : les distributions de longueurs, la répétitivité des formats et
la diversité des sorties diffèrent. Une faible perte peut notamment provenir
d'un livrable très structuré et prévisible. Elle ne prouve ni que l'outil choisi
est pertinent, ni que son argument est valide, ni que le constat final est
correct. La perte plus élevée de la reconnaissance peut, inversement, refléter
une plus grande diversité de plans et d'observations sans établir une moindre
efficacité opérationnelle. Ces résultats restent donc des contrôles de
convergence, pas une validation du \gls{hmoe} agentique.

\section{Protocole d'évaluation et ablations}
\label{sec:hmoe-evaluation}

La validation de l'hypothèse de la \cref{sec:hmoe-probleme} doit comparer des
systèmes soumis aux mêmes scénarios, prompts, outils, budgets de génération et
graines. Le jeu d'évaluation doit être distinct des traces d'entraînement et
inclure des topologies ou variantes non observées. Quatre conditions isolent
les apports : le modèle de base avec prompt seul ; un adaptateur QLoRA global
entraîné sur le corpus réuni ; les quatre experts avec sélection forcée de
l'expert correct ; enfin le \gls{hmoe} avec routage automatique. Une ablation
supplémentaire, qui permute volontairement certains experts, mesure si le gain
provient bien de la spécialisation plutôt que du seul ajout de paramètres.

L'évaluation doit être conduite sur deux niveaux complémentaires :

\begin{description}
  \item[Conformité agentique] taux d'appels d'outils syntaxiquement valides,
    respect des outils autorisés, achèvement des phases, conformité des
    livrables et fréquence des récupérations après erreur ;
  \item[Efficacité de l'audit] précision, rappel et score F1 des constats,
    proportion de preuves vérifiées, couverture des chemins d'attaque et coût
    en jetons, temps et mémoire.
\end{description}

Les décisions sémantiques peuvent être assistées par un juge \gls{llm}, mais
celui-ci ne doit pas être l'unique arbitre : les biais de position, de verbosité
et d'auto-préférence sont documentés \cite{zheng2023judge}. Les validations de
schéma, la correspondance avec la \gls{veriteterrain} et le calcul des métriques
doivent rester déterministes lorsque cela est possible. Il faudra enfin publier
les intervalles de confiance sur plusieurs exécutions et journaliser la version
du modèle de base, des adaptateurs, du routeur, des corpus et du code. Ce
protocole permettra de conclure séparément sur trois questions : QLoRA
apporte-t-il un gain par rapport au prompt seul, la séparation des experts
apporte-t-elle un gain par rapport à un adaptateur global, et le routeur
automatique conserve-t-il ce gain sans connaître à l'avance l'expert correct ?

\chapter{Acquisition de compétences d'ingénieur}
\label{chap:competences}

\textbf{Structure provisoire :} le référentiel officiel de compétences de
TELECOM Nancy n'a pas été retrouvé dans les espaces de travail. Les situations
ci-dessous sont factuelles, mais leurs intitulés et niveaux devront être alignés
sur ce référentiel avant la remise.

\section{Concevoir un système d'intelligence artificielle sous contraintes}

La réalisation du \gls{hmoe} m'a conduit à traiter conjointement apprentissage,
architecture logicielle et exploitation. La contrainte matérielle interdisait
de maintenir quatre modèles complets. J'ai donc partagé une base quantifiée et
spécialisé son comportement avec quatre adaptateurs \gls{qlora}. Ce choix ne
suffisait cependant pas à obtenir un agent fiable : il a fallu préparer des
conversations conservant les couples appel--résultat, aligner les fenêtres
d'entraînement et d'inférence, sérialiser le changement d'adaptateur et limiter
la surface d'outils présentée à chaque expert.

Cette situation m'a appris à évaluer un choix de modèle à l'échelle du système.
La consommation mémoire, la qualité des données, le routage, les contrats
d'outils et le comportement en cas de sortie malformée sont interdépendants.
Une bonne perte de validation ne remplace donc ni les tests d'intégration ni une
campagne agentique sur des scénarios représentatifs.

\section{Analyser et maîtriser un système complexe}

LANCE relie topologie réseau, outils protocolaires, orchestration multi-phase,
modèle de langage et infrastructure de test. Pour éviter qu'une erreur se
propage silencieusement, j'ai transformé plusieurs attentes implicites en
invariants vérifiables : prérequis de phase, couverture minimale de
reconnaissance, provenance d'une preuve, périmètre réseau autorisé et filtrage
des constats avant le rapport final.

\section{Expérimenter et valider}

Le harness d'évaluation m'a obligé à distinguer une affirmation plausible d'un
résultat démontré. J'ai travaillé sur une correspondance structurelle entre
constats et vérité terrain, sur la recomposition du niveau de preuve depuis les
appels d'outils, puis sur la couverture de chemins ordonnés. Les erreurs
d'outils, les tests négatifs concluants et l'absence de test reçoivent des
statuts différents. Cette distinction évite qu'une panne soit interprétée comme
une absence de vulnérabilité ou qu'un texte assuré soit crédité sans exécution.

J'ai également appris à séparer les indicateurs internes des conclusions
scientifiques. Les pertes d'entraînement servent à détecter une divergence ou à
choisir un checkpoint ; elles ne mesurent pas directement la capacité d'un
expert à terminer une phase. De même, l'instantané expérimental actuel constitue
surtout un test d'intégration du pipeline et du harness. Une comparaison
contrôlée du \gls{hmoe} avec plusieurs baselines reste nécessaire.

\section{Industrialiser et documenter}

Le passage du prototype au service a demandé une frontière explicite entre code
versionné, données, checkpoints et adaptateurs finaux. La synchronisation repose
sur une liste blanche, des empreintes et des limites de taille ; le serveur est
exposé par une \gls{api} compatible avec le reste de LANCE et supervisé comme un
service. Côté benchmark, la composition des scénarios, leur déploiement et leur
évaluation partagent des contrats versionnés. J'ai ainsi appris qu'une
fonctionnalité n'est réellement exploitable que si elle peut être déployée,
observée, testée et reproduite.

\section{Auto-évaluation}

Je considère avoir progressé dans la conception d'architectures hybrides, où un
modèle probabiliste reste encadré par des mécanismes déterministes. Je dois
encore approfondir deux points : construire une campagne d'ablation permettant
d'attribuer précisément les gains au routage des experts, et calibrer les
heuristiques du graphe sur des observations empiriques. Le niveau officiel de
chaque compétence restera à compléter à partir du référentiel de TELECOM Nancy.

\chapter{Conclusion}
\label{chap:conclusion}

Ce stage avait pour objectif principal de spécialiser LANCE par une architecture
HMoE agentique, avec une empreinte d'entraînement compatible avec des ressources
limitées. La solution partage un modèle de base quantifié entre plusieurs
adaptateurs QLoRA associés aux rôles de la campagne. Un routeur déterministe
sélectionne l'expert attendu selon la phase et le serveur expose l'ensemble
derrière une interface commune. La mission a couvert la préparation des traces,
l'entraînement, le chargement des adaptateurs et leur intégration au cycle
d'action de l'agent.

Le pipeline et le harness ont été consolidés pour rendre cette spécialisation
évaluable. Le pipeline encadre les transitions, restreint les outils et valide
les livrables. Le harness déploie des scénarios reproductibles, associe les
résultats à une vérité terrain et mesure la qualité des campagnes. Les travaux
sur l'API, le suivi des exécutions, l'automatisation de l'infrastructure et les
métriques participent donc à un même protocole expérimental. Le graphe pondéré
et le cône d'impact apportent en complément un signal de priorité lorsque l'état
du réseau évolue.

Les résultats provisoires montrent que l'agent peut achever des campagnes sur
plusieurs scénarios, mais aussi que la production d'un constat ne garantit pas
un chemin d'attaque entièrement prouvé. Cet écart confirme l'intérêt d'évaluer
le modèle au niveau de son exécution plutôt qu'à partir de la seule
vraisemblance de ses réponses.

Plusieurs limites demeurent. Les adaptateurs dépendent de la diversité et de la
correction des traces. Le routage par phase est interprétable, mais ne combine
pas plusieurs spécialités aux frontières des rôles. Les poids du graphe restent
heuristiques et les premiers résultats ne constituent pas encore une comparaison
définitive avec un modèle généraliste.

La suite devra exécuter les configurations sur un jeu de scénarios identique,
répéter les campagnes et analyser leur qualité, leurs preuves et leurs coûts.
Une boucle d'amélioration pourra ensuite soumettre les erreurs à une revue
humaine avant d'enrichir les données. Ce projet m'a finalement appris qu'un
modèle spécialisé ne peut être jugé indépendamment de l'architecture qui
encadre ses actions et du protocole qui rend ses résultats vérifiables.


% ─── Bibliographie ───────────────────────────────────────────────────────────
\clearpage
\begin{singlespace}
\printbibliography[heading=bibintoc, title={Bibliographie}]
\end{singlespace}

% ─── Annexes ─────────────────────────────────────────────────────────────────
% N'active les annexes que si elles sont réellement appelées dans le corps du texte.
% \appendix
% \chapter{<<Titre de l'annexe A>>}

% ⚠️ Une annexe doit être APPELÉE dans le corps du texte (\cref{annexe:xxx}),
% sinon elle ne compte pas — et la plupart des écoles le vérifient.
% Contenu typique : extraits de code longs, schémas grand format, protocoles
% de test détaillés, captures d'écran de l'interface.


\end{document}
